% =========================================================
%  全国大学生数学建模竞赛论文模板
%  文档类：cumcmthesis（需安装或 Overleaf 导入）
%
%  编译方式：xelatex（必须）
%    xelatex cumcm2026_template.tex   % 第 1 次
%    xelatex cumcm2026_template.tex   % 第 2 次（消交叉引用警告）
%  推荐：latexmk -xelatex cumcm2026_template.tex
%
%  2025 年电子版提交要求：
%    - 去掉封面与编号页（加 withoutpreface 选项）
%    - 不需要目录
%    - 不要承诺书
% =========================================================

% 为 xcolor 启用 table 选项（\rowcolor）
\PassOptionsToPackage{table}{xcolor}
\documentclass[withoutpreface,bwprint]{cumcmthesis}

% 打印版（含封面，需手动贴承诺书）
% \documentclass{cumcmthesis}

% ─── 额外宏包（cumcmthesis 已内置大多数常用宏包）───────────
\usepackage{bookmark}
\usepackage{listings}

% ─── listings 全局样式（Python 语法高亮）───────────────────
\lstset{
    language=Python,
    basicstyle=\ttfamily\small,
    keywordstyle=\color{blue!70!black},
    commentstyle=\color{green!50!black},
    stringstyle=\color{red!60!black},
    numbers=left, numberstyle=\tiny\color{gray},
    numbersep=8pt, breaklines=true,
    columns=flexible, aboveskip=0pt, belowskip=0pt,
    frame=none, tabsize=4,
}

% ─── longtable 三线表（可跨页，水平居中）──────────────
\setlength{\LTleft}{0pt plus 1fill}
\setlength{\LTright}{0pt plus 1fill}

% ─── 公式间距统一 ───────────────────────────────────────
% 公式-文字、文字-文字、公式-公式的间距保持一致：
% 全部等于一行行距（\baselineskip）。如需微调，改这里即可。
\setlength{\abovedisplayskip}{\baselineskip}
\setlength{\belowdisplayskip}{\baselineskip}
\setlength{\abovedisplayshortskip}{\baselineskip}
\setlength{\belowdisplayshortskip}{\baselineskip}

% ─── 排版约定 ──────────────────────────────────────────
% 1. \cref{} 前不要加"图/表/假设"前缀（\cref 自动带标签词，会重复）
% 2. 中文引号一律用 “”（U+201C/U+201D），不用英式直引号 "（编译后同向）
% 3. 符号说明直接三线表、无说明文字，且不加 \caption（不参与表编号）
% 4. 每张三线表加 \setlength{\tabcolsep}{...}（默认 6pt，方便调宽度）

% ─── \scite 命令（CUMCM 上标引用格式）─────────────────────
\newcommand{\scite}[1]{\textsuperscript{[#1]}}

% ─── 论文基本信息 ──────────────────────────────────────────
\title{【在此填写论文题目】}
\tihao{A}                      % 题号：A / B / C / D / E
\baominghao{0000}              % 报名号（12 位）
\schoolname{XX大学}
\membera{成员甲}
\memberb{成员乙}
\memberc{成员丙}
\supervisor{指导教师}
\yearinput{2025}
\monthinput{09}
\dayinput{07}

% =========================================================
\begin{document}

\maketitle

% ─── 摘要 ─────────────────────────────────────────────────
\begin{abstract}

    【简述问题背景与现实意义。1~2 句点明要解决的核心问题。】

    针对问题一，建立了【模型名】模型，采用【算法】求解，得到【量化结果与对比基准】。

    针对问题二，基于问题一的结果，进一步建立了【模型名】模型，通过【方法】分析，得到【量化结果】。

    针对问题三，综合考虑【新增因素】，建立了【模型名】模型，利用【算法】求解，得到【量化结果】。

    灵敏度分析表明，模型对关键参数具有良好的鲁棒性。【补充灵敏度具体数值。】

    \keywords{数学模型 \quad XXX模型 \quad YYY方法 \quad ZZZ算法}
\end{abstract}

% 目录（2019 年起明确不要目录，保留注释供参考）
% \tableofcontents

% =========================================================
\section{问题重述}
% =========================================================

\subsection{问题背景}

【描述题目背景，1~2 段，简明扼要，不要大段复制题目原文。说明该问题的现实意义和研究的必要性。】

\subsection{问题提出}

根据题目要求，需要解决以下问题：

\begin{enumerate}
    \item \textbf{问题一}：【简述问题一的核心任务，体现自己的理解。】
    \item \textbf{问题二}：【在问题一基础上，新增了什么条件或要求。】
    \item \textbf{问题三}：【进一步拓展或深化的问题。】
\end{enumerate}

% =========================================================
\section{问题分析}
% =========================================================

\subsection{问题一的分析}

【针对问题一的特点，分析需要用什么思路来建模。说明关键变量、约束条件、数据特征，以及为什么选择某种方法。】

\subsection{问题二的分析}

【在问题一的基础上，分析新增因素带来的建模挑战和改进方向。】

\subsection{问题三的分析}

【分析问题的复杂性，说明需要引入什么新方法或视角。每段末尾点明所用模型的核心思想。】

% =========================================================
\section{模型假设}
% =========================================================

为使模型易于建立与分析，在不失一般性的前提下，作出以下假设：

\begin{assumption}
    【假设一】假设……，理由是……
\end{assumption}

\begin{assumption}
    【假设二】假设……，理由是……
\end{assumption}

\begin{assumption}
    【假设三】假设……，理由是……
\end{assumption}

% =========================================================
\section{符号说明}
% =========================================================

% 符号说明：直接三线表、无说明文字；不加 \caption（不参与表编号）
{\renewcommand{\arraystretch}{1.38}
\setlength{\tabcolsep}{6pt}
\begin{longtable}{ccl}
    \toprule[1.5pt]
    符号       & 单位 & 说明     \\
    \midrule[1pt]
    \endfirsthead
    \endhead
    \bottomrule[1.5pt]
    \endlastfoot
    $x$      & —  & 【变量说明】 \\
    $t$      & s  & 时间     \\
    $\alpha$ & —  & 【参数说明】 \\
    $f(x)$   & —  & 目标函数   \\
\end{longtable}}

% =========================================================
\section{模型准备}
% =========================================================

\subsection{数据整理与处理}

【说明数据来源、格式、规模。缺失值处理、异常值识别、归一化/标准化方法。】

\begin{figure}[H]
    \centering
    % \includegraphics[width=0.7\textwidth]{figures/data_distribution.pdf}
    \fbox{\parbox{0.7\textwidth}{\centering\vspace{2cm}【在此插入数据分布图】\vspace{2cm}}}
    \caption{数据分布情况}
    \label{fig:data_dist}
\end{figure}

\subsection{可视化分析}

【对数据进行初步可视化，分析数据特征（分布、相关性、趋势等），为模型选择提供依据。展示关键发现。】
如 \cref{fig:correlation} 所示，【变量A】与【变量B】呈现【正/负】相关，说明……

\begin{figure}[H]
    \centering
    % \includegraphics[width=0.6\textwidth]{figures/correlation.pdf}
    \fbox{\parbox{0.6\textwidth}{\centering\vspace{2cm}【在此插入相关性分析图】\vspace{2cm}}}
    \caption{相关性分析}
    \label{fig:correlation}
\end{figure}

\subsection{技术路线}

本文的整体建模流程如 \cref{fig:flowchart} 所示。

\begin{figure}[H]
    \centering
    % \includegraphics[width=0.85\textwidth]{figures/flowchart.pdf}
    \fbox{\parbox{0.85\textwidth}{\centering\vspace{3cm}【在此插入技术路线图/流程图】\vspace{3cm}}}
    \caption{技术路线图}
    \label{fig:flowchart}
\end{figure}

% =========================================================
\section{问题一的模型建立及求解}
% =========================================================

\subsection{模型建立}

【阐述建模思路。首先定义变量，然后建立目标函数和约束条件。相关研究见文献\scite{1}。】

设【变量定义】……，则问题可转化为如下优化模型：

\begin{equation}
    \min_{x \in \mathbb{R}^n} \quad f(x) = \sum_{i=1}^{n} w_i \cdot g_i(x)
    \label{eq:obj1}
\end{equation}

约束条件为：

\begin{equation}
    \text{s.t.} \quad
    \begin{cases}
        h_1(x) \leq 0 \\
        h_2(x) = 0    \\
        x \geq 0
    \end{cases}
    \label{eq:con1}
\end{equation}

其中，\cref{eq:obj1} 表示……，$w_i$ 为权重系数，$g_i(x)$ 为第 $i$ 项指标函数。
\cref{eq:con1} 中各约束分别表示……

\subsection{模型求解}

采用【算法名】对上述模型进行求解。步骤如下：

\begin{enumerate}
    \item 步骤一：【数据处理/初始化】
    \item 步骤二：【核心迭代/寻优过程】
    \item 步骤三：【结果输出/后处理】
\end{enumerate}

【可在此处插入求解流程图，算法实现代码见附录。】

\subsection{结果分析}

求解结果如表 \cref{tab:result1} 所示，图 \cref{fig:result1} 展示了【关键对比/趋势】。

{\renewcommand{\arraystretch}{1.38}
\setlength{\tabcolsep}{6pt}
\begin{longtable}{ccccc}
    \caption{问题一计算结果}\label{tab:result1}\\
    \toprule[1.5pt]
    参数       & 方案 1 & 方案 2 & 方案 3 & 最优  \\
    \midrule[1pt]
    \endfirsthead
    \endhead
    \bottomrule[1.5pt]
    \endlastfoot
    $\alpha$ & 0.1  & 0.3  & 0.5  & 0.3 \\
    目标值      & 12.5 & 8.7  & 10.2 & 8.7 \\
    耗时 (s)   & 1.2  & 2.3  & 1.8  & 2.3 \\
\end{longtable}}

\begin{figure}[H]
    \centering
    % \includegraphics[width=0.7\textwidth]{figures/result1.pdf}
    \fbox{\parbox{0.7\textwidth}{\centering\vspace{3cm}【在此插入结果图】\vspace{3cm}}}
    \caption{问题一结果图示}
    \label{fig:result1}
\end{figure}

由表 \cref{tab:result1} 可以看出……。图 \cref{fig:result1} 表明……
【量化分析结果的意义，与基准方案对比，给出具体改善幅度。】

% =========================================================
\section{问题二的模型建立及求解}
% =========================================================

\subsection{模型建立}

在问题一的基础上，考虑【新增因素】，对模型进行扩展。

\begin{equation}
    J = \int_0^T f\!\left(x(t),\, u(t)\right)\,\mathrm{d}t
    \label{eq:obj2}
\end{equation}

【详细定义变量和约束……】

\subsection{模型求解}

采用【算法/方法】进行求解。核心思路是……

\begin{enumerate}
    \item 【第一步】
    \item 【第二步】
    \item 【第三步】
\end{enumerate}

【求解实现代码见附录。】

\subsection{结果分析}

求解结果如表 \cref{tab:result2} 所示。

{\renewcommand{\arraystretch}{1.38}
\setlength{\tabcolsep}{6pt}
\begin{longtable}{ccc}
    \caption{问题二计算结果}\label{tab:result2}\\
    \toprule[1.5pt]
    指标  & 基准方案 & 优化方案       \\
    \midrule[1pt]
    \endfirsthead
    \endhead
    \bottomrule[1.5pt]
    \endlastfoot
    指标A & XX   & YY（改善ZZ\%） \\
    指标B & XX   & YY（改善ZZ\%） \\
\end{longtable}}

【对结果进行量化分析，说明优化效果和原因。】

% =========================================================
\section{问题三的模型建立及求解}
% =========================================================

\subsection{模型建立}

【问题三的建模描述，说明与前两问的区别和联系。】

\begin{equation}
    \max_{x,y} \quad F(x,y) = \lambda_1 f_1(x) + \lambda_2 f_2(y)
    \label{eq:obj3}
\end{equation}

满足约束：
\begin{equation}
    \text{s.t.} \quad
    \begin{cases}
        g_1(x,y) \leq 0 \\
        g_2(x) \leq C   \\
        y \in \{0,1\}
    \end{cases}
    \label{eq:con3}
\end{equation}

其中 $\lambda_1 + \lambda_2 = 1$ 为权重系数……

\subsection{模型求解}

【说明求解方法，如果是多目标优化则说明权重设定或 Pareto 方法。】

\subsection{结果分析}

求解结果如表 \cref{tab:result3} 所示。

{\renewcommand{\arraystretch}{1.38}
\setlength{\tabcolsep}{6pt}
\begin{longtable}{lccc}
    \caption{问题三对比结果}\label{tab:result3}\\
    \toprule[1.5pt]
    方案   & 效率指标 & 成本指标 & 综合评价 \\
    \midrule[1pt]
    \endfirsthead
    \endhead
    \bottomrule[1.5pt]
    \endlastfoot
    方案A  & X1   & Y1   & Z1   \\
    方案B  & X2   & Y2   & Z2   \\
    推荐方案 & X*   & Y*   & Z*   \\
\end{longtable}}

\begin{figure}[H]
    \centering
    % \includegraphics[width=0.7\textwidth]{figures/result3.pdf}
    \fbox{\parbox{0.7\textwidth}{\centering\vspace{3cm}【在此插入多目标对比图/Pareto前沿图】\vspace{3cm}}}
    \caption{多目标优化结果对比}
    \label{fig:result3}
\end{figure}

【分析结果，解释为什么推荐方案最优，给出量化依据。】

% =========================================================
\section{模型测试与分析}
% =========================================================

\subsection{参数灵敏度分析}

为检验模型对关键参数的稳定性，对参数 $\alpha$ 进行灵敏度分析。
保持其他参数不变，令 $\alpha$ 在基准值 $\alpha_0 = 0.3$ 的
$\pm 20\%$ 范围内变化，观察目标函数值的变化，结果如图 \cref{fig:sensitivity} 所示。

\begin{figure}[H]
    \centering
    % \includegraphics[width=0.65\textwidth]{figures/sensitivity.pdf}
    \fbox{\parbox{0.65\textwidth}{\centering\vspace{3cm}【在此插入灵敏度分析图】\vspace{3cm}}}
    \caption{参数灵敏度分析}
    \label{fig:sensitivity}
\end{figure}

由图可知，当 $\alpha$ 在 $[0.24,\, 0.36]$ 范围内变化时，
目标函数值变化幅度不超过 $X\%$，说明模型对该参数具有较好的鲁棒性。

\subsection{权重灵敏度分析}

【若模型涉及主观权重（如 AHP、TOPSIS、加权多目标），需同时做权重灵敏度分析。】

当权重 $w$ 在 $[0.3, 0.7]$ 范围内变化时，最优方案排序保持稳定/发生以下变化……
【说明结果是否对权重的选择敏感。】

\subsection{模型稳健性分析}

【可选：通过情景构造（如需求扰动、能力下降、参数组合变动）进一步验证模型的适应性。】

% =========================================================
\section{模型评价与总结}
% =========================================================

\subsection{模型的优点}

\begin{enumerate}
    \item \textbf{【优点一】}：【具体说明，尽量量化。】
    \item \textbf{【优点二】}：【如：模型融合了XXX和YYY的优点，兼顾了解释性和精度。】
    \item \textbf{【优点三】}：【如：灵敏度分析表明结论在参数合理范围内保持稳健。】
\end{enumerate}

\subsection{模型的改进}

\begin{enumerate}
    \item \textbf{【改进一】}：【当前模型的局限性是什么，未来如何改进。】
    \item \textbf{【改进二】}：【如：当前假设简化了XXX，未来可考虑XXX。】
\end{enumerate}

\subsection{模型的推广}

【说明模型可以迁移到哪些其他领域或问题类型。】

% =========================================================
%  AI 工具使用声明（2026 年试行规定，2026-09-01 起施行）
%  ◇ 必须位于参考文献【之前】，二选一，本队固定为已使用版本（见下）
%  ◇ 支撑材料必须附《AI工具使用详情.pdf》（见 ai-usage-detail-template.tex）
%     ——缺此文件按违规处理，故意隐瞒/虚假声明可被取消评奖资格
%  ◇ 若确未使用 AI，替换为：
%  \section*{AI工具使用声明}
%  本参赛队在竞赛过程中未使用任何AI工具。
%  正文声明与支撑材料详情需保持一致
\CumcmAIUsed{语言润色、代码调试、文献整理}
% =========================================================
%  参考文献
%  引用方式：\scite{编号} 如 \scite{1}、\scite{2,3}、\scite{4-6}
% =========================================================
\begin{thebibliography}{99}

    \bibitem{ref-book}
    作者名. \textit{文献标题}\allowbreak[M]. 出版社, 出版地, 年份.

    \bibitem{ref-journal}
    Author A, Author B. Title of the article\allowbreak[J]. \textit{Journal Name}, Year, Vol(Issue): Pages.

    \bibitem{ref-standard}
    全国大学生数学建模竞赛论文格式规范（2023 年修订）\allowbreak[S]. 全国大学生数学建模竞赛组委会, 2023.

    \bibitem{ref-url}
    作者名. 网页标题\allowbreak[EB/OL]. \url{https://example.com}, 访问日期.

\end{thebibliography}

\newpage
% =========================================================
\section*{附录}
% =========================================================
% ─── 将代码写入外部文件（供 \lstinputlisting 读取）───────
\begin{filecontents*}[overwrite]{appendix1.py}
import numpy as np
import pandas as pd

# ==========================
# 1. 基础数据输入
# ==========================
# 突水点涌水量 (m^3/min)
Q = 30.0
# 巷道宽度 (m)
W = 4.0
# 巷道高度 (m)
H = 3.0

def calc_flow_time(length, flow_rate, area=0.4):
    """计算巷道过水耗时"""
    volume = area * length
    return volume / flow_rate

# 单巷道过水时间计算
tunnel_length = 150.0
t = calc_flow_time(tunnel_length, Q)
print(f"巷道过水耗时: {t:.2f} min")
\end{filecontents*}
\begin{filecontents*}[overwrite]{appendix2.py}
import heapq

def dynamic_dijkstra(graph, start, end, water_time):
    """动态 Dijkstra（考虑水位变化）"""
    dist = {node: float('inf') for node in graph}
    prev = {node: None for node in graph}
    dist[start] = 0.0
    pq = [(0.0, start)]

    while pq:
        d, u = heapq.heappop(pq)
        if u == end:
            break
        for v, length, tid in graph[u]:
            if tid in water_time:
                speed = 2.0 if d >= water_time[tid] else 4.0
            else:
                speed = 4.0
            nd = d + length / speed
            if nd < dist[v]:
                dist[v] = nd
                prev[v] = u
                heapq.heappush(pq, (nd, v))

    path, node = [], end
    while node is not None:
        path.append(node)
        node = prev[node]
    path.reverse()
    return dist[end], path
\end{filecontents*}
\begin{filecontents*}[overwrite]{appendix3.py}
import matplotlib.pyplot as plt
import numpy as np

COLOR_BLUE = '#1f4e79'
COLOR_ORANGE = '#ed7d31'

def plot_sensitivity(param_name, param_values, obj_values):
    """绘制参数灵敏度分析曲线"""
    plt.figure(figsize=(8, 5))
    plt.plot(param_values, obj_values,
             color=COLOR_BLUE, marker='o', linewidth=2)
    plt.axhline(y=obj_values[0], color='gray',
                linestyle='--', alpha=0.7, label='基准值')
    plt.xlabel(param_name)
    plt.ylabel('目标函数值')
    plt.title(f'{param_name} 灵敏度分析')
    plt.legend()
    plt.grid(True, alpha=0.3)
    plt.tight_layout()
    plt.savefig('fig_sensitivity.pdf', dpi=300)
\end{filecontents*}

\begin{appendices}
    \section{支撑材料文件列表}

    本论文支撑材料清单如下，全部文件随论文一并提交。

    {\renewcommand{\arraystretch}{1.38}
    \setlength{\tabcolsep}{6pt}
    \begin{longtable}{cllc}
        \caption{支撑材料文件列表}\\
        \toprule[1.5pt]
        序号 & 文件名 & 说明 & 类型 \\
        \midrule[1pt]
        \endfirsthead
        \endhead
        \bottomrule[1.5pt]
        \endlastfoot
        1 & appendix1.py & 问题一求解程序 & 源程序 \\
        2 & appendix2.py & 问题二求解程序 & 源程序 \\
        3 & appendix3.py & 问题三求解程序 & 源程序 \\
        4 & 【data/xxx.csv】 & 【原始数据】 & 数据文件 \\
    \end{longtable}}

    \section{源程序代码}

    本部分给出建模所用到的全部完整、可运行的源程序代码。

    \subsection{第一问求解程序}
    \lstinputlisting[frame=tb, framerule=0.8pt, rulecolor=\color{gray!50}]{appendix1.py}

    \subsection{第二问求解程序}
    \lstinputlisting[frame=tb, framerule=0.8pt, rulecolor=\color{gray!50}]{appendix2.py}

    \subsection{第三问求解程序}
    \lstinputlisting[frame=tb, framerule=0.8pt, rulecolor=\color{gray!50}]{appendix3.py}

    \section{补充数据}

    【在此放原始数据表格或中间计算结果。】

\end{appendices}

\end{document}
